\paragraph{What the diagnostics show.} Four findings recur across our
experiments. First, the supracompetitive outcomes of Q-learning informed
traders do not rest on punishment: rivals do not react to deviations, even
with perfect monitoring or when deviations are hundreds of noise standard
deviations large, and deviating pays. Second, the same learning procedure
under-trades with no rival at all, by an amount that grows with
$\sqrt\alpha$ and disappears under counterfactual updates; in the market,
counterfactual updating restores competition while degrading information does
nothing. Third, placebos that cannot collude reproduce the conventional
evidence of collusion: myopic learners score at or beyond the cartel on
intensity- and informativeness-based indices, an uninformative memory
reproduces the effect of an informative one, and a myopic shared-table learner
shows the threshold-shaped reaction to price shocks that has been read as a
price trigger. Fourth, detectability rather than learning sophistication
decides whether punishment is even feasible: in the standard Kyle market a
deviation is a fraction of a noise standard deviation
(Proposition~\ref{prop:snr}).

\paragraph{Implications for auditing algorithms.} Section~\ref{sec:audit}
turns the diagnostics into a protocol. Its value is that every step can fail:
applied to the \citet{calvano2020} baseline it finds punishment, and applied
to informed traders it does not.

\paragraph{Implications for policy.} If supracompetitive trading is a learning
bias, interventions aimed at collusion through information, such as delayed or
coarser order-flow disclosure, will not work; our reduced-transparency
treatment confirms this. What works is algorithm design: counterfactual
(full-information) updating removes the bias. This mirrors conclusions of
\citet{asker2022,asker2024} and \citet{abada2023} for pricing algorithms and
extends them to informed trading. A bias is not harmless: it still lowers
liquidity and price informativeness, myopic learners push trading below even
the cartel level, and with passive traders it costs the learners themselves
\PassiveProfitPct\% of their equilibrium profit. But its remedy is a
design standard, not a cartel investigation.

\paragraph{Limitations.} Our evidence concerns tabular Q-learning with standard
hyperparameters and, more briefly, DQN and PPO; learners that model their
rivals may find punishment strategies these do not. Constant-step Q-learning
does not converge in the strict sense of \citet{calvano2020}: greedy
strategies still change on path by one to two grid steps over the last
$10^5$ periods, although outcomes are stable across seeds, schedules and
horizons. Our main training horizons ($6\times10^6$ to $1.5\times10^7$
periods) are shorter than those of \citet{dou2025} and \citet{esquinas2026}
(up to $5\times10^{10}$ and $5\times10^8$). Runs of $2\times10^8$ periods
in the $\xi=500$ calibration reproduce the separate-table index of
\citet{esquinas2026} and leave every diagnostic unchanged, and runs with the
stopping rule of \citet{dou2025}, which every session meets within
$2\times10^9$ periods, change nothing. We cannot rule out that other
implementation details let their learners discover punishment, which is why
the placebo, not the absence of a response, carries our argument there. Even at $2\times10^8$ periods greedy strategies keep changing in most
on-path states ($\LongSharedPolicyChange$ of entries over the last $10^7$
periods for the shared table), as expected with a constant step size. The market maker learns its price impact from
exponentially weighted moments rather than being fixed at the equilibrium
value; replacing it by the rolling least-squares market maker of
\citet{dou2025} leaves the results unchanged. Finally, the stylised market abstracts from inventory dynamics,
multiple assets and long-lived information.

\paragraph{Conclusion.} Low trading by learning algorithms is easy to produce
and easy to misread. Telling collusion from learning bias requires
diagnostics that can fail: a deviation test, placebos that cannot collude, and
a check of whether the outcome moves with the learning rule. In the canonical
market for informed trading, and in the regime where deviations are most
visible, these diagnostics point to pruning, not punishment.
